\chapter{实践方法与验收}

\section{原理}

操作系统教材的实践目标不是写一个玩具内核名字，而是把关键状态机写清楚。调度器保护的是任务状态；页表保护的是地址和权限；文件描述符表保护的是对象生命周期。每个模块都必须能输出诊断，让错误路径可复查。

验收的顺序是：状态机正确，权限边界正确，生命周期正确，等待和关闭正确，持久化边界正确，然后再讨论性能和源码阅读。没有这些基础，读 Linux 源码也只会变成函数名堆砌。

本册暂时不展开实践卷工程，但理论卷也必须有实践方法。每章都要能落成一份小报告：问题是什么，抽象对象是什么，不变量是什么，正常路径如何推进，错误路径如何退出，哪些证据证明结论，哪些真实内核行为没有覆盖。

源码阅读要服务现象。不要从内核源码第一个文件开始顺序读，也不要把函数名搬进笔记。正确方式是从一个现象入口进入：为什么线程睡眠、为什么缺页、为什么 fd 仍有效、为什么 \code{write} 可能短写、为什么 \code{fsync} 慢、为什么信号打断系统调用。读源码时只追踪足以解释现象的对象和状态。

验收不是追求“模型覆盖真实 OS 全部细节”。教学模型必须小，但小模型不能破坏核心契约。调度模型可以没有 CFS/EEVDF 细节，但不能调度 blocked 线程；页表模型可以没有多级页表，但不能忽略权限；fd 模型可以没有完整 VFS，但不能把 dup 后引用关系写错；持久化模型可以不实现真实文件系统，但不能把 write 成功当成业务提交。

\section{实践}

理论卷推荐使用下面的报告结构：

\begin{lstlisting}
Title:
  要解释的 OS 现象或机制

State:
  涉及的内核对象、用户态对象、硬件对象

Invariant:
  必须始终成立的关系

Operation:
  正常路径的状态转移

Failure:
  错误路径、取消路径、关闭路径、崩溃点

Evidence:
  命令、trace、日志、计数器、测试断言

Boundary:
  当前模型没有覆盖的真实内核细节
\end{lstlisting}

每章至少完成一个纸上状态机。调度章写 runqueue；同步章写有界队列；锁章写锁顺序图；系统调用章写参数检查；exec 章写提交点；内存章写页表翻译和缺页；mmap 章写 VMA 到 page cache 的路径；块层章写 bio 到 request 的转换；I/O 章写 fd 引用图；socket 章写 accept/send/recv 等待队列；panic 章写失败回滚；持久化章写提交状态机；观测章写一次性能或卡死诊断报告。

实践时要避免两个极端。第一个极端是只画概念图，不写错误路径；第二个极端是直接读真实内核大段源码，却没有抽象状态。合格学习路径是先有小状态机，再用真实系统现象校验它的边界。

\section{逐章项目}

本册理论卷可以配一条纸上项目线，不需要先写完整内核：

\begin{enumerate}
\tightlist
\item 写一个裸机超级循环，列出每个任务的周期和最坏执行时间。
\item 把一个阻塞 I/O 函数改成状态机，证明它每次 step 都会返回。
\item 加入中断事件队列，证明 ISR 只做短路径工作。
\item 写启动状态机：BSS、data、早期 console、页表、调度器、init 任务。
\item 加入定时器和时间片，写出上下文保存字段。
\item 实现 FCFS、RR、优先级和 MLFQ 的调度表模拟。
\item 写 lost wakeup 测试，证明等待队列和谓词由同一把锁保护。
\item 为有界队列写等待谓词、关闭路径和背压策略。
\item 写锁顺序图、lost wakeup 反例和 lockdep 检查思路。
\item 写系统调用分发、用户指针复制和错误码矩阵。
\item 写 exec 装载流程：ELF 验证、段映射、用户栈、CLOEXEC、提交点。
\item 手算页表翻译、COW 写故障和 Clock 页面置换。
\item 画 mmap 缺页路径，区分 shared/private、dirty page 和 msync。
\item 实现 buddy/slab 的纸上分配释放轨迹，检查碎片和守恒。
\item 画驱动描述符环、DMA 映射、completion 和 reset 清理路径。
\item 设计块层 bio 合并、deadline 调度、flush/FUA 顺序测试。
\item 画 fd 引用图、page cache 读写路径和异步 completion。
\item 设计 inode、目录项、bitmap 和 VFS 的最小文件系统模型。
\item 画 socket bind/listen/accept/send/recv 的状态机和等待队列。
\item 设计临时文件、fsync、rename、manifest 的可靠提交协议。
\item 写 capability、namespace、cgroup、seccomp 的最小权限报告。
\item 写 errno、信号、WARN、panic 的错误分类和回滚路径报告。
\item 为一次故障写现象、假设、证据、反证、边界报告。
\end{enumerate}

\section{评分矩阵}

\begin{longtable}{p{0.22\linewidth}p{0.34\linewidth}p{0.32\linewidth}}
\toprule
能力 & 合格 & 高质量 \\
\midrule
状态机 & 能列出状态和转换 & 能列出错误、取消、关闭、崩溃边界 \\
算法 & 能写伪代码 & 能分析复杂度、参数和不变量 \\
实现 & 能说明数据结构 & 能说明并发、权限和生命周期风险 \\
测试 & 覆盖正常路径 & 覆盖边界路径、反例和故障注入 \\
观测 & 有日志或命令输出 & 有可复现证据、反证和未验证边界 \\
\bottomrule
\end{longtable}

\section{测试}

最低验收包括：round-robin 调度、blocked/terminated 过滤、锁顺序检查、页表权限、unmap、exec 失败回滚、fd dup、fd close 和未知句柄诊断。扩展验收包括：优先级 aging、copy-on-write、mmap shared/private、块层 flush 顺序、socket 非阻塞等待、短写、崩溃恢复和设备阻塞。

\begin{rubricbox}
\begin{itemize}
\tightlist
\item 调度：runnable、running、blocked、zombie、reaped 的状态转换可解释。
\item 同步：等待谓词完整，关闭路径能唤醒所有等待者。
\item 锁：锁保护对象明确，锁顺序无环，持锁上下文不违反睡眠规则。
\item 内存：地址翻译保留 offset，权限错误和合法缺页能区分。
\item syscall：用户指针、fd、权限和错误码都有检查。
\item exec：ELF 验证、段映射、用户栈和提交点可解释。
\item mmap：VMA、PTE、page cache 与 dirty/writeback 状态一致。
\item I/O：fd、open file description、inode、page cache 和设备请求不混淆。
\item 驱动：描述符环、DMA buffer 和中断下半部生命周期正确。
\item 块层：bio 合并、request 派发、flush/FUA/barrier 顺序可验证。
\item 文件系统：路径解析、inode、目录项和空闲空间管理可解释。
\item socket：阻塞、非阻塞、poll、epoll 与等待队列关系明确。
\item 安全：最小权限、namespace、cgroup、seccomp 有明确边界。
\item 错误恢复：errno、信号、WARN、panic 和回滚路径不混用。
\item 持久化：write、fsync、rename、manifest 的提交边界清楚。
\item 观测：每个结论有证据和反证，不把猜测写成事实。
\item 源码阅读：能说明入口、对象、状态变化、返回给用户态的现象。
\end{itemize}
\end{rubricbox}

测试失败时要先问是哪条不变量断了：任务状态、地址权限、对象引用、等待谓词、提交边界还是恢复规则。把失败归到不变量，比只看报错信息更接近真实 OS 调试。
